\section{Removing an unknown watermark: literature and experiment
selection}\label{record-docs-watermark-removal-literature-md-removing-an-unknown-watermark-literature-and-experiment-selection}

Reviewed on 2026-10-26. The target is a model whose weights and ordinary
architecture are accessible, while its watermark technology, key,
message and extractor remain unknown. An API-only generator is a
different access setting.

The shortlist contains \textbf{two model-level method families}. They
require no watermark detector queries. These are research candidates,
not established universal removers.

{\scriptsize\setlength{\tabcolsep}{3pt}\def\LTcaptype{none} % do not increment counter
\begin{longtable}[]{@{}
  >{\raggedright\arraybackslash}p{(\linewidth - 4\tabcolsep) * \real{0.3333}}
  >{\raggedright\arraybackslash}p{(\linewidth - 4\tabcolsep) * \real{0.3333}}
  >{\raggedright\arraybackslash}p{(\linewidth - 4\tabcolsep) * \real{0.3333}}@{}}
\toprule\noalign{}
\begin{minipage}[b]{\linewidth}\raggedright
Method
\end{minipage} & \begin{minipage}[b]{\linewidth}\raggedright
Clean weights available
\end{minipage} & \begin{minipage}[b]{\linewidth}\raggedright
Clean weights unavailable
\end{minipage} \\
\midrule\noalign{}
\endhead
\bottomrule\noalign{}
\endlastfoot
Merge toward the compatible clean checkpoint & Primary inexpensive
attack; use a small fixed interpolation sweep & Ineligible; do not
quietly substitute a related checkpoint \\
Continued training on independent clean data & Eligible without using
the clean weights; permits comparison with merging & Primary attack,
starting only from the watermarked weights \\
\end{longtable}
}

For merging, full replacement is a sanity control. The meaningful
question is whether partial merging removes detectable signals while
retaining useful capabilities. A parameter difference between two
checkpoints does not identify which changes encode a watermark.

For continued training, use the architecture's ordinary objective:
next-token prediction for an LLM or autoregressive image backbone, and
denoising for a diffusion backbone. Broad representative data are the
primary choice. Existing LoRA runs provide a cheaper baseline;
full-parameter training is the stronger candidate where resources
permit. These are variants within one family, not an additional
collection of attacks. The trainable modules must be chosen from the
ordinary training recipe, without consulting the actual watermark
location. Failure against a watermark in another component must remain a
reported failure.

\subsection{Evidence from full
papers}\label{record-docs-watermark-removal-literature-md-evidence-from-full-papers}

Full PDFs, including appendices, were downloaded and read. Local
research assets are under
\texttt{\seqsplit{results/research\_workflow/watermarking/resources/}}.

\begin{itemize}
\item
  \textbf{\href{https://proceedings.mlr.press/v306/rebuffi26a.html}{Learning
  to Watermark in the Latent Space of Generative Models}, ICML 2026.}
  The peer-reviewed DistSeal version directly studies clean-data LoRA
  forgetting in Appendix G. After 2,500 steps on 50 images of one class,
  the strongest reported learning rate reduces mean bit accuracy to
  roughly 0.70 for diffusion and 0.65 for RAR. This supports testing
  clean training on our image backbones. It establishes weakening, not
  complete removal: bit accuracy is not the detection rate at a fixed
  false-positive rate. The experiment also does not establish
  broad-domain removal or preserved quality after every attack.
\item
  \textbf{\href{https://www.sri.inf.ethz.ch/publications/scharff26mergewm}{Making
  Open-Source Text LLM Watermarks Durable Against Merging}, NeurIPS 2026
  listing.} The authors' public 26-page PDF reports substantial removal
  by merging with clean weights, including standard AAR/KTH distilled
  watermarks. Crucially, their merge-adversarial training preserves much
  more signal under the same operations. Merging is therefore a strong
  candidate with explicit counterexamples to universality. Appendix A.2
  also shows ordinary fine-tuning substantially weakens detection, with
  domain-dependent effects. The title is listed in the conference's
  \href{https://neurips.cc/Downloads/2026}{2026 downloads}; the
  downloadable author PDF is still labeled ``Preprint.''
\item
  \textbf{\href{https://www.sri.inf.ethz.ch/publications/gloaguen2025oss}{Towards
  Watermarking of Open-Source LLMs}, WMARK workshop at ICLR 2025.} This
  is workshop evidence, not an ICLR main-conference paper. Its
  systematic comparison supports merging and broad-domain clean
  training. Full training weakens the tested watermarks more
  consistently than LoRA; task-specific training can leave substantial
  signal outside the training domain. Ordinary quantization and pruning
  often preserve the watermark. This motivates excluding a broad
  compression/pruning grid from the initial shortlist.
\item
  \textbf{\href{https://proceedings.iclr.cc/paper_files/paper/2025/hash/c0f7ee1901fef1da4dae2b88dfd43195-Abstract-Conference.html}{Black-Box
  Detection of Language Model Watermarks}, ICLR 2025.} Prompt-based
  statistical tests detect evidence of three watermark families without
  the private key, a clean counterpart or detector queries. They test
  context-dependent preferences, repeated-output diversity and cache
  effects. Their scope is those mechanisms and their modeling
  assumptions; they do not establish a detector for arbitrary unseen
  watermark technologies. Their instruction-following setup and
  family-specific assumptions need validation before applying them to
  instruct Qwen, TextSeal 2 or our learned latent adaptation.
\item
  \textbf{\href{https://doi.org/10.1007/978-981-95-3537-8_30}{When There
  Is No Decoder}, ICICS 2025.} The full author version is
  arXiv:2507.03646. Despite its no-box setting, its fine-tuning attack
  trains a surrogate watermark extractor and an alternative payload.
  Results depend on assumptions about the extractor's capacity; some
  image-quality comparisons use only 20 images. It is less suitable as
  our first scheme-agnostic method than ordinary clean training, and is
  not selected.
\end{itemize}

\subsection{Detection, selection and
verification}\label{record-docs-watermark-removal-literature-md-detection-selection-and-verification}

Knowing that a model is watermarked does not require first recovering
its technology before attempting removal. Without an owner detector, a
negative generic test cannot certify that an unknown watermark is
absent. A classifier that separates real from generated content is not a
watermark detector.

With a clean counterpart, paired-output differences or a learned
classifier can measure distribution changes, but those changes can
include ordinary fine-tuning effects. Without the clean counterpart,
that paired comparison is unavailable. The existing base/student text
classifier is therefore an evaluator diagnostic, not an attacker
capability in the no-clean-weights setup.

The implementation uses separate inputs and outputs:

\begin{enumerate}
\def\labelenumi{\arabic{enumi}.}
\tightlist
\item
  The attacker receives the watermarked model, architecture/tokenizer,
  its independent clean fitting data, a fixed compute budget and utility
  checks. Only the first setup additionally receives clean model
  weights. It receives no watermark identifier, implementation, key,
  payload, extractor or score.
\item
  Attack strengths and checkpoints are fixed beforehand or selected
  using clean validation utility. Owner detection results cannot select
  the attack, its hyperparameters, updated modules or stopping point.
\item
  A separate evaluator applies the actual owner detector after artifacts
  are frozen, using held-out inputs and a frozen threshold. Report
  detection, false positives with uncertainty, quality and compute,
  including performance outside the fitting domain. Suppression of a
  measured signal is the supported claim; universal erasure is not.
\end{enumerate}

The original merge and clean-training jobs are useful baselines. Their
training losses do not use watermark scores, but training and owner
evaluation currently share orchestration and some stages expose fixed
intermediate measurements. They should not be described as an enforced
blind access boundary. The new
\texttt{\seqsplit{scripts/watermark\_removal.py}} workflow reuses their
training loops with separate validated worker contracts and a later
owner-evaluation process. It preserves the original jobs and their
frozen snapshots. The boundary controls inputs and feedback in this
research protocol; it is not a security sandbox for arbitrary code
running under the same filesystem account.

TextSeal is a separate sampler-level control: its Qwen weights are
unchanged. Weight merging alone does not remove a watermark reapplied by
an external sampler. Ordinary local inference tests sampler bypass; it
is not evidence of removing an unknown in-model watermark. Per-output
rewriting or image regeneration would likewise clean individual outputs
rather than establish model-weight removal, so they are outside this
initial two-method shortlist.

The current text experiment supersedes the initial shortlist for text.
See the \href{records/docs/instruct_watermarking.tex}{instruction-model
protocol} for TextSeal 2 and the DistSeal-inspired student, BIRA, direct
resampling and smoothing. Earlier Base-model watermark outputs and
metrics were deleted; image results remain.
